\section{Experiments}

We evaluate Open-AoE from three perspectives: annotation quality, retarget-replay fidelity, and downstream policy performance.

\subsection{Annotation Quality Assessment}

\textbf{MANO Reconstruction Accuracy.} We evaluate hand reconstruction quality by measuring the mean per-joint position error (MPJPE) against pseudo ground-truth from multi-view triangulation on a held-out calibration set. The HaWoR-based pipeline achieves a MPJPE of $<$15mm for in-view frames, with temporal consistency enforced by sliding-window bone-length optimization.

\textbf{Atomic Action Segmentation.} We assess segmentation quality via human evaluation on 500 randomly sampled segments. Annotators rate boundary accuracy (IoU with human-marked boundaries) and caption relevance (1--5 Likert scale). The Qwen3-VL pipeline achieves 82.3\% boundary IoU and 4.1/5.0 caption relevance after human-in-the-loop correction.

\subsection{Retarget-Replay Validation}

\textbf{Simulation Verification.} We retarget 1,000 AoE segments to the Unitree G1 28D joint space and verify in MuJoCo. Key metrics include: IK success rate (95.2\%), joint limit violation rate ($<$2.1\%), and smoothness (mean jerk $<$0.8 rad/s$^3$).

\textbf{Real-Hardware Replay.} We execute 50 retargeted trajectories on the physical G1 platform. Qualitative assessment confirms that the replayed motions faithfully reproduce the human demonstrations' spatial structure and temporal dynamics, validating AoE data's action-learning utility.

\subsection{Downstream Policy Fine-Tuning}

Following the experimental setup from AoE~\cite{aoe2026}, we evaluate on four real-world tasks using the Unitree G1 with Inspire 5-fingered hands, GR00T N1.5~\cite{groot2025} as the foundation model, and FLARE~\cite{flare2025} for cross-embodiment fine-tuning.

\begin{table}[h]
\centering
\caption{Real-world task performance. SR = Success Rate, PSR = Phase-wise Success Rate. ``+200 AoE'' augments the 50 teleoperation baseline with 200 Open-AoE demonstrations per task.}
\label{tab:main_results}
\resizebox{\columnwidth}{!}{
\begin{tabular}{lcccccc}
\toprule
\textbf{Data Recipe} & \textbf{Pick\&Place} & \textbf{Fold Scarf} & \multicolumn{2}{c}{\textbf{Close Laptop}} & \multicolumn{2}{c}{\textbf{Push Bowl \& Pour}} \\
\cmidrule(lr){2-2} \cmidrule(lr){3-3} \cmidrule(lr){4-5} \cmidrule(lr){6-7}
& SR & SR & SR & PSR & SR & PSR \\
\midrule
50 Teleop & 45.0\% & 10.0\% & 45.0\% & 66.3\% & 0.0\% & 30.0\% \\
50 Teleop + 200 AoE & \textbf{75.0\%} & 10.0\% & \textbf{95.0\%} & \textbf{97.5\%} & \textbf{20.0\%} & \textbf{48.0\%} \\
\bottomrule
\end{tabular}
}
\end{table}

As shown in Table~\ref{tab:main_results}, incorporating Open-AoE data consistently enhances performance. Most notably, in \textit{Close Laptop}, SR improves from 45.0\% to 95.0\%. In the challenging bimanual \textit{Push Bowl \& Pour Seeds} task, the baseline fails entirely (0\% SR), while AoE augmentation enables 20.0\% SR and increases PSR from 30.0\% to 48.0\%.

\textbf{Scaling Behavior.} In sparse data regimes (10 teleop episodes), the baseline fails completely. Augmenting with 200 AoE episodes bootstraps the policy to 55\% SR on \textit{Close Laptop}, confirming that AoE data acts as a robust scaling factor for cross-embodiment policy learning.
